\documentclass[UTF8,12pt,fontset=windows]{ctexart}
\usepackage[paperwidth=240mm,paperheight=200mm,left=14mm,right=14mm,top=9mm,bottom=12mm]{geometry}
\usepackage{amsmath,amssymb,mathtools,bm,graphicx,xcolor}
\pagestyle{empty}
\setlength{\parindent}{0pt}
\setlength{\parskip}{4pt}
\setlength{\emergencystretch}{2em}
\xeCJKDeclareCharClass{CJK}{"2460 -> "24FF}
\xeCJKDeclareCharClass{CJK}{"2160 -> "2188}
\begin{document}
{\large\bfseries 2026.1}

设 $x\to0$ 时，函数 $f(x)=ax^2+bx+\arcsin x$ 和 $g(x)=\sqrt[3]{1+x^2}-1$ 是等价无穷小，则（　　）

(A) $a=\dfrac13$，$b=-1$.

(B) $a=-\dfrac13$，$b=-1$.

(C) $a=\dfrac13$，$b=1$.

(D) $a=-\dfrac13$，$b=1$.

\vfill
\newpage
{\large\bfseries 2026.2}

设 $y_1(x),y_2(x)$ 是某非齐次线性微分方程的两个特解，若常数 $\lambda,\mu$ 使得 $2\lambda y_1(x)+\mu y_2(x)$ 是该方程的解，$\lambda y_1(x)-2\mu y_2$ 是该方程对应的齐次方程的解，则（　　）

(A) $\lambda=\dfrac15,\mu=\dfrac25$.

(B) $\lambda=\dfrac25,\mu=\dfrac15$.

(C) $\lambda=\dfrac14,\mu=\dfrac12$.

(D) $\lambda=\dfrac12,\mu=\dfrac14$.

\vfill
\newpage
{\large\bfseries 2026.3}

设函数 $z=z(x,y)$ 由方程 $x-az=\mathrm{e}^{y+az}$（$a$ 是非零常数）确定，则（　　）

(A) $\dfrac{\partial z}{\partial x}-\dfrac{\partial z}{\partial y}=\dfrac1a$.

(B) $\dfrac{\partial z}{\partial x}+\dfrac{\partial z}{\partial y}=\dfrac1a$.

(C) $\dfrac{\partial z}{\partial x}-\dfrac{\partial z}{\partial y}=-\dfrac1a$.

(D) $\dfrac{\partial z}{\partial x}+\dfrac{\partial z}{\partial y}=-\dfrac1a$.

\vfill
\newpage
{\large\bfseries 2026.4}

设线密度为 $1$ 的均匀细棒的两个端点分别为 $(-1,0)$，$(1,0)$，在 $(0,1)$ 处有一个质量为 $m$ 的质点，$G$ 为引力常数，则细棒对质点的万有引力为（　　）

(A) $2\displaystyle\int_0^1\dfrac{Gmx}{(x^2+1)^{\frac12}}\,\mathrm{d}x$.

(B) $2\displaystyle\int_0^1\dfrac{Gm}{(x^2+1)^{\frac12}}\,\mathrm{d}x$.

(C) $2\displaystyle\int_0^1\dfrac{Gmx}{(x^2+1)^{\frac32}}\,\mathrm{d}x$.

(D) $2\displaystyle\int_0^1\dfrac{Gm}{(x^2+1)^{\frac32}}\,\mathrm{d}x$.

\vfill
\newpage
{\large\bfseries 2026.5}

设函数 $f(x)$ 在区间 $[-1,1]$ 上有定义，则（　　）

(A) 当 $f(x)$ 在 $(-1,0)$ 单调递减，在 $(0,1)$ 单调递增时，$f(0)$ 是极小值.

(B) 当 $f(0)$ 是极小值时，$f(x)$ 在 $(-1,0)$ 单调递减，在 $(0,1)$ 单调递增.

(C) 当 $f(x)$ 的图形在 $[-1,1]$ 是凹的时，$\dfrac{f(x)-f(1)}{x-1}$ 在 $[-1,1)$ 单调递增.

(D) 当 $\dfrac{f(x)-f(1)}{x-1}$ 在 $[-1,1)$ 单调递增时，$f(x)$ 的图形在 $[-1,1]$ 是凹的.

\vfill
\newpage
{\large\bfseries 2026.6}

$f(x)=\displaystyle\int_1^{x^3}\dfrac{\mathrm{e}^t}{1+t^2}\,\mathrm{d}t$，$f(x)$ 的反函数为 $g(x)$，则（　　）

(A) $g(0)=1$，$g'(0)=\dfrac32\mathrm{e}$.

(B) $g(0)=1$，$g'(0)=\dfrac{2}{3\mathrm{e}}$.

(C) $g(1)=0$，$g'(1)=\dfrac32\mathrm{e}$.

(D) $g(1)=1$，$g'(1)=\dfrac{2}{3\mathrm{e}}$.

\vfill
\newpage
{\large\bfseries 2026.7}

设函数 $f(x,y)$ 在平面区域 $D=\{(x,y)\mid 0\leqslant x\leqslant1,\;0\leqslant y\leqslant1\}$ 上连续，$f(x,y)=f(y,x)$，则 $\displaystyle\iint_D f(x,y)\,\mathrm{d}x\mathrm{d}y=$（　　）

(A) $2\lim\limits_{n\to\infty}\sum_{i=1}^{n}\sum_{j=n+1-i}^{n}f\left(\dfrac{i}{n},\dfrac{j}{n}\right)\dfrac{1}{n^2}$.

(B) $\dfrac12\lim\limits_{n\to\infty}\sum_{i=1}^{n}\sum_{j=1}^{n}f\left(\dfrac{i}{n},\dfrac{j}{n}\right)\dfrac{1}{n^2}$.

(C) $2\lim\limits_{n\to\infty}\sum_{i=1}^{2n}\sum_{j=1}^{2n+1-i}f\left(\dfrac{i}{2n},\dfrac{j}{2n}\right)\dfrac{1}{n^2}$.

(D) $\dfrac12\lim\limits_{n\to\infty}\sum_{i=1}^{2n}\sum_{j=1}^{i}f\left(\dfrac{i}{2n},\dfrac{j}{2n}\right)\dfrac{1}{n^2}$.

\vfill
\newpage
{\large\bfseries 2026.8}

单位矩阵经过若干次互换两行得到的矩阵称为置换矩阵。设 $A$ 为 $n$ 阶置换矩阵，$A^*$ 为 $A$ 的伴随矩阵，则（　　）

(A) $A^*$ 为置换矩阵.

(B) $A^{-1}$ 为置换矩阵.

(C) $A^{-1}=A^*$.

(D) $A^{-1}=-A^*$.

\vfill
\newpage
{\large\bfseries 2026.9}

$A=\begin{pmatrix}1&0&1\\0&0&1\\1&1&3\\1&1&1\end{pmatrix}$，$C=\begin{pmatrix}2&0\\1&1\\1&1\\a&b\end{pmatrix}$，存在矩阵 $B$ 满足 $AB=C$，则（　　）

(A) $a=-1$，$b=-1$.

(B) $a=2$，$b=2$.

(C) $a=-1$，$b=2$.

(D) $a=2$，$b=-1$.

\vfill
\newpage
{\large\bfseries 2026.10}

设 $3$ 阶矩阵 $A,B$ 满足 $AB+BA=A^2+B^2$ 且 $A\neq B$，以下结论错误的是（　　）

(A) $(A-B)^3=O$.

(B) $A-B$ 只有零特征值.

(C) $A,B$ 不能都是对角矩阵.

(D) $A-B$ 只有一个线性无关的特征向量.

\vfill
\newpage
{\large\bfseries 2026.11}

$p$ 为常数，反常积分 $\displaystyle\int_0^{+\infty}\dfrac{\arctan x}{x^p(x+1)}\,\mathrm{d}x$ 收敛，则 $p$ 的取值范围是\underline{\hspace{20mm}}.

\vfill
\newpage
{\large\bfseries 2026.12}

$\lim\limits_{x\to0}\left(\dfrac1x-\dfrac{\ln(1+x)}{x\sin x}\right)=$\underline{\hspace{20mm}}.

\vfill
\newpage
{\large\bfseries 2026.13}

曲线 $x^2+2\sqrt3\,xy+y^2=1$ 在 $(0,1)$ 处的曲率半径为\underline{\hspace{20mm}}.

\vfill
\newpage
{\large\bfseries 2026.14}

设函数 $f(x,y)$ 可微，且满足 $\mathrm{d}f(0,0)=\pi\,\mathrm{d}x+3\,\mathrm{d}y$，$g(x)=f(\ln x,\sin\pi x)$，则 $g'(1)=$\underline{\hspace{20mm}}.

\vfill
\newpage
{\large\bfseries 2026.15}

函数 $f(x)=\ln(2+x)$ 在区间 $[0,2]$ 上的平均值为\underline{\hspace{20mm}}.

\vfill
\newpage
{\large\bfseries 2026.16}

矩阵 $A=\begin{pmatrix}1&b&-1\\a+2&3&-3\end{pmatrix}$，二次型 $x^{\mathrm{T}}(AA^{\mathrm{T}})x$ 的规范形为 $y_1^2$，则 $a+b=$\underline{\hspace{20mm}}.

\vfill
\newpage
{\large\bfseries 2026.17}

（本题满分10分）

求积分 $I=\displaystyle\int_{-1}^{1}\mathrm{d}x\int_{|x|}^{\sqrt{2-x^2}}y\sin\sqrt{x^2+y^2}\,\mathrm{d}y$.

\vfill
\newpage
{\large\bfseries 2026.18}

（本题满分12分）

已知 $g(x)$ 为连续函数，$f(x)=\displaystyle\int_0^{x^2}g(xt)\,\mathrm{d}t$，求 $f'(x)$，且判断 $f'(x)$ 在 $x=0$ 处的连续性.

\vfill
\newpage
{\large\bfseries 2026.19}

（本题满分12分）

求函数 $f(x,y)=(2x^2-y^2)\mathrm{e}^x$ 的极值.

\vfill
\newpage
{\large\bfseries 2026.20}

（本题满分12分）

设 $M(x_0,y_0)$ 为曲线 $y=\dfrac{1}{1+x^2}\ (x\geqslant0)$ 的拐点. $O$ 为坐标原点，设平面区域 $D$ 是由第一象限内的曲线 $y=\dfrac{1}{1+x^2}\ (x\geqslant x_0)$、线段 $OM$ 以及 $x$ 轴正半轴围成的无界区域. 求 $D$ 绕 $x$ 轴所成的旋转体的体积.

\vfill
\newpage
{\large\bfseries 2026.21}

（本题满分12分）

求微分方程 $x^2y''-2xy'-(y')^2=0$ 满足 $y\Big|_{x=3}=\dfrac12$，$y'\Big|_{x=3}=-9$ 的特解.

\vfill
\newpage
{\large\bfseries 2026.22}

（本题满分12分）

向量 $\boldsymbol{\alpha}_1=\begin{pmatrix}1\\0\\-1\\-1\end{pmatrix}$，$\boldsymbol{\alpha}_2=\begin{pmatrix}1\\-1\\0\\-2\end{pmatrix}$，$\boldsymbol{\alpha}_3=\begin{pmatrix}0\\-1\\1\\-1\end{pmatrix}$，$\boldsymbol{\alpha}_4=\begin{pmatrix}0\\1\\-1\\1\end{pmatrix}$. 记 $A=(\boldsymbol{\alpha}_1,\boldsymbol{\alpha}_2,\boldsymbol{\alpha}_3,\boldsymbol{\alpha}_4)$，$G=(\boldsymbol{\alpha}_1,\boldsymbol{\alpha}_2)$.

（1）证明：$\boldsymbol{\alpha}_1,\boldsymbol{\alpha}_2$ 是 $\boldsymbol{\alpha}_1,\boldsymbol{\alpha}_2,\boldsymbol{\alpha}_3,\boldsymbol{\alpha}_4$ 的极大线性无关组；

（2）求矩阵 $H$ 使得 $A=GH$，并求 $A^{10}$.

\vfill
\end{document}
